\documentclass[a4paper,12pt]{article}
\usepackage[utf8]{inputenc}
\usepackage[T1]{fontenc}
\usepackage[spanish]{babel}
\usepackage{geometry}
\usepackage{graphicx}
\usepackage{xcolor}
\usepackage{hyperref}
\usepackage{fancyhdr}
\usepackage{enumitem}
\usepackage{amsmath}
\usepackage{booktabs}
\usepackage{array}
\usepackage{setspace}
\usepackage{float}
\usepackage{longtable}
\usepackage{tabularx}
\usepackage{ragged2e}

\geometry{top=34mm,bottom=25mm,left=25mm,right=25mm}

\definecolor{ExternadoGreen}{RGB}{0,104,55}
\definecolor{ExternadoDark}{RGB}{0,51,25}
\definecolor{ExternadoGold}{RGB}{198,146,20}

\hypersetup{colorlinks=true,linkcolor=ExternadoGreen,urlcolor=ExternadoGreen,citecolor=ExternadoGreen}

\setlength{\headheight}{52pt}
\pagestyle{fancy}
\fancyhf{}
\fancyhead[L]{\IfFileExists{firma.jpg}{\includegraphics[height=40pt]{firma.jpg}}{}}
\fancyhead[R]{\textcolor{ExternadoGreen}{\small Departamento de Matemáticas}}
\fancyfoot[C]{\textcolor{ExternadoDark}{\thepage}}

\fancypagestyle{plain}{
  \fancyhf{}
  \fancyhead[L]{\IfFileExists{firma.jpg}{\includegraphics[height=40pt]{firma.jpg}}{}}
  \fancyhead[R]{\textcolor{ExternadoGreen}{\small Universidad Externado de Colombia}}
  \fancyfoot[C]{\textcolor{ExternadoDark}{\thepage}}
}

\title{
\vspace{-8mm}
\textbf{\textcolor{ExternadoGreen}{Ficha de Proyecto de Investigación}}\\
\large \textcolor{ExternadoDark}{Grupo de Investigación en Ciencia de Datos}\\
\large \textcolor{ExternadoDark}{Universidad Externado de Colombia}
}
\author{\textcolor{ExternadoDark}{Departamento de Matemáticas}}
\date{\textcolor{ExternadoDark}{\today}}

\setlist[itemize]{leftmargin=*, itemsep=3pt, topsep=3pt}
\setstretch{1.12}

\begin{document}
\maketitle
\vspace{-3mm}
\noindent\textcolor{ExternadoDark}{\rule{\textwidth}{0.6pt}}

\section*{1. Información general}
\begin{longtable}{p{4.2cm} p{10.4cm}}
\toprule
\textbf{Campo} & \textbf{Detalle} \\
\midrule
Título del proyecto & Optimización de algoritmos de ruta \\
Investigador principal & David Andrés Franco Quintero \\
Co-investigadores & No reportado \\
Líneas de investigación & Modelado Matemático, Machine Learning \\
Año de inicio & 2024-04-01 \\
Duración estimada & No reportado. \\
Estado del proyecto & Formulación \\
Tipo de proyecto & Consultoría / transferencia; \\
Nivel de avance actual & Idea inicial; \\
Semillero vinculado & No \\
Nombre del semillero & No aplica \\
Fuentes de financiación & No requiere financiación; \\
\bottomrule
\end{longtable}

\section*{2. Problema de investigación}
¿Cómo rediseñar el esquema de rutas de Transmilenio teniendo presente la inclusión de nuevos factores como la Línea 1 del Metro y la troncal Av. 68, con el fin de generar eficiencias en el sistema?

\section*{3. Objetivo general}
Especificar un conjunto nuevo de rutas que deben implementarse en el sistema Transmilenio a fin de optimizar los trayectos hechos por los usuarios.

\section*{4. Metodología}
Alcance del proyecto: proposición de esquema de rutas nuevas en Transmilenio. Necesidad de información: datos de validación de usuarios, datos de rutas que paran actualmente, datos de estaciones actuales y futuras. Estas son obtenidas de los datos públicos de Transmilenio. Fuentes de información: Datos Abiertos de Transmilenio. Estrategias de acceso de datos: bases institucionales. Aliados externos: escucha de la propuesta de cambio de rutas, mostrando la eficiencia generada con ello.

\section*{5. Comunidades a impactar}
\subsection*{5.1. Comunidades externas}
\begin{itemize}
    \item Población usuaria de Transmilenio Ciudadano común que usa Transporte Público Sistema Integrado de Transporte Público de Bogotá
\end{itemize}

\subsection*{5.2. Comunidades internas}
No reportado.

\section*{6. Semillero y articulación formativa}
\subsection*{6.1. Participación del semillero}
\begin{itemize}
    \item Recolección de datos
    \item Desarrollo metodológico
    \item Programación / implementación
    \item Análisis de resultados
    \item Redacción de productos
\end{itemize}

\subsection*{6.2. Aliados externos}
No reportado.

\section*{7. Requerimientos tecnológicos}
\begin{itemize}
    \item Herramientas de visualización
    \item Computación de alto rendimiento (HPC)
\end{itemize}

\section*{8. Resultados esperados}
\subsection*{8.1. Resultados generales}
- Análisis actual del componente Transmilenio en términos de tiempo de desplazamiento. - Análisis del componente Transmilenio (en términos de tiempo de desplazamiento) con las rutas actuales y asumiendo una entrada de la Linea 1 del Metro y la troncal Av. 68. - - Análisis del componente Transmilenio (en términos de tiempo de desplazamiento) con un nuevo sistema de rutas y asumiendo una entrada de la Linea 1 del Metro y la troncal Av. 68.

\subsection*{8.2. Resultados científicos esperados}
\begin{itemize}
    \item Artículo científico
\end{itemize}

\subsection*{8.3. Productos aplicados}
\begin{itemize}
    \item Ruteo del componente Transmilenio
\end{itemize}

\section*{9. Observación de gestión}
Este documento fue generado automáticamente a partir del formulario de registro de proyectos del grupo. Se recomienda revisar y ajustar la redacción final antes de su circulación institucional o envío a convocatorias.

\end{document}
